\documentclass[11pt]{article}

\usepackage[margin=1in]{geometry}
\usepackage{amsmath}
\usepackage{amssymb}
\usepackage{mathtools}
\usepackage{enumitem}
\usepackage{xcolor}

\newcommand{\R}{\mathbb{R}}
\newcommand{\ve}[1]{\boldsymbol{#1}}
\newcommand{\tr}{\intercal}
\DeclareMathOperator{\diag}{diag}
\DeclareMathOperator{\Trace}{Tr}

\title{Kalman Associative Memory Derivations}
\author{}
\date{}

\begin{document}
\maketitle

\section{Kalman Associative Memory: optimal update}

Let the associative-memory state be
\[
  \ve{S}_t \in \R^{d_k\times d_v},
  \qquad
  \mathcal{M}_t(\ve{k})=\ve{S}_t^{\tr}\ve{k}.
\]
The linear--Gaussian state-space model is
\begin{align}
  \ve{S}_t^\star
  &= \ve{D}_t\ve{S}_{t-1}^\star+\ve{W}_t,
  &
  \ve{W}_t&\sim\mathcal{N}(0,\ve{\Omega}_t), \\
  \ve{v}_t
  &= \ve{S}_t^{\star\tr}\ve{k}_t+\ve{e}_t,
  &
  \ve{e}_t&\sim\mathcal{N}(0,r_t\ve{I}).
\end{align}
Let $\ve{z}_s=(\ve{k}_s,\ve{v}_s)$ denote the key--value observation at step $s$.
At the end of step $t-1$, the Kalman filter represents the posterior over the latent
memory as a Gaussian,
\begin{equation}
  p(\ve{S}_{t-1}^\star\mid \ve{z}_{1:t-1})
  =
  \mathcal{N}(\ve{S}_{t-1},\ve{P}_{t-1}),
  \label{eq:previous-filtering-posterior}
\end{equation}
where $\ve{S}_{t-1}$ is the posterior mean, i.e.\ the point estimate of the memory after
processing the first $t-1$ key--value pairs, and $\ve{P}_{t-1}$ is its key-space uncertainty.
For the matrix memory, this covariance is understood per value coordinate, with the same
key-space covariance shared across value columns.

Before observing token $t$, the filter first predicts the next latent memory by marginalizing
over the previous posterior:
\begin{equation}
  p(\ve{S}_t^\star\mid \ve{z}_{1:t-1})
  =
  \int
  p(\ve{S}_t^\star\mid \ve{S}_{t-1}^\star)
  p(\ve{S}_{t-1}^\star\mid \ve{z}_{1:t-1})
  \,d\ve{S}_{t-1}^\star .
  \label{eq:prediction-integral}
\end{equation}
Because both factors are linear--Gaussian, the predicted distribution is Gaussian,
\begin{equation}
  p(\ve{S}_t^\star\mid \ve{z}_{1:t-1})
  =
  \mathcal{N}(\widehat{\ve{S}}_t,\widehat{\ve{P}}_t).
  \label{eq:predicted-prior}
\end{equation}
Its mean and covariance are the exact Kalman prediction:
\begin{equation}
  \widehat{\ve{S}}_t=\ve{D}_t\ve{S}_{t-1},
  \qquad
  \widehat{\ve{P}}_t
  =
  \ve{D}_t\ve{P}_{t-1}\ve{D}_t^{\tr}+\ve{\Omega}_t.
  \label{eq:kam-predict}
\end{equation}
Thus $\widehat{\ve{S}}_t$ is a point estimate, but specifically the \emph{predicted prior
mean} of $\ve{S}_t^\star$ conditioned on the past key--value sequence
$\ve{z}_{1:t-1}$. It is not yet the point estimate of
$p(\ve{S}_t^\star\mid \ve{z}_{1:t})$, because the current pair
$(\ve{k}_t,\ve{v}_t)$ has not been used. The covariance $\widehat{\ve{P}}_t$ is not a point
estimate; it is the predicted uncertainty around that mean.
The predicted read and innovation are
\begin{equation}
  \widehat{\ve{v}}_t=\widehat{\ve{S}}_t^{\tr}\ve{k}_t,
  \qquad
  \ve{\delta}_t=\ve{v}_t-\widehat{\ve{v}}_t.
  \label{eq:kam-innovation}
\end{equation}

This residual form is the matrix version of the standard Kalman correction. To see it,
first look at one value coordinate $j$. Let $\ve{s}_{t,j}$ be column $j$ of the memory and
let $v_{t,j}$ be coordinate $j$ of the observed value. The measurement model is
\begin{equation}
  v_{t,j}
  =
  \ve{k}_t^{\tr}\ve{s}_{t,j}^\star+e_{t,j},
  \qquad
  e_{t,j}\sim\mathcal{N}(0,r_t).
  \label{eq:single-value-coordinate-observation}
\end{equation}
Before this measurement is used, the predicted distribution for this column is
\begin{equation}
  \ve{s}_{t,j}^\star \mid \ve{z}_{1:t-1}
  \sim
  \mathcal{N}(\widehat{\ve{s}}_{t,j},\widehat{\ve{P}}_t).
  \label{eq:single-coordinate-prior}
\end{equation}
The observation $v_{t,j}$ is one scalar. Its predicted mean is
$\ve{k}_t^{\tr}\widehat{\ve{s}}_{t,j}$, so the part of the observation not already predicted
by the prior is the scalar residual
\begin{equation}
  \delta_{t,j}
  =
  v_{t,j}-\ve{k}_t^{\tr}\widehat{\ve{s}}_{t,j}.
  \label{eq:single-coordinate-innovation}
\end{equation}
For a joint Gaussian pair $(x,y)$ with scalar $y$, the conditional mean is
\[
  \mathbb{E}[x\mid y]
  =
  \mathbb{E}[x]
  +
  \operatorname{Cov}(x,y)\operatorname{Var}(y)^{-1}
  \big(y-\mathbb{E}[y]\big).
\]
Apply this identity with $x=\ve{s}_{t,j}^\star$ and $y=v_{t,j}$. Under
\eqref{eq:single-coordinate-prior}--\eqref{eq:single-value-coordinate-observation},
\begin{equation}
  \operatorname{Cov}(\ve{s}_{t,j}^\star,v_{t,j})
  =
  \widehat{\ve{P}}_t\ve{k}_t,
  \qquad
  \operatorname{Var}(v_{t,j})
  =
  \ve{k}_t^{\tr}\widehat{\ve{P}}_t\ve{k}_t+r_t.
  \label{eq:single-coordinate-cross-cov}
\end{equation}
Therefore the posterior mean is
\begin{equation}
  \ve{s}_{t,j}
  =
  \widehat{\ve{s}}_{t,j}
  +
  \ve{\kappa}_t\,\delta_{t,j}.
  \label{eq:single-coordinate-kalman-correction}
\end{equation}
where
\begin{equation}
  \ve{\kappa}_t
  =
  \frac{\widehat{\ve{P}}_t\ve{k}_t}
  {\ve{k}_t^{\tr}\widehat{\ve{P}}_t\ve{k}_t+r_t}.
  \label{eq:single-coordinate-gain}
\end{equation}
The same $\ve{\kappa}_t$ applies to every value coordinate because the model uses the same
key-space covariance $\widehat{\ve{P}}_t$ and the same observation noise $r_t$ for every
coordinate of $\ve{v}_t$. Stacking \eqref{eq:single-coordinate-kalman-correction} over
$j=1,\ldots,d_v$ gives the rank-one residual write
\begin{equation}
  \ve{S}_t
  =
  \widehat{\ve{S}}_t
  +
  \ve{\kappa}_t
  \big(\ve{v}_t-\widehat{\ve{S}}_t^{\tr}\ve{k}_t\big)^{\tr}.
  \label{eq:kam-residual-write}
\end{equation}
Thus the residual write does not come from DeltaNet first; it comes from the Kalman
posterior-mean update $x^+=x^-+K(y-Hx^-)$. DeltaNet later appears as the fixed-gain special
case where $\ve{\kappa}_t$ is forced to be proportional to $\ve{k}_t$.
For a candidate write direction $\ve{\kappa}_t$, the posterior covariance induced by this
write is
\begin{equation}
  \ve{P}_t(\ve{\kappa}_t)
  =
  (\ve{I}-\ve{\kappa}_t\ve{k}_t^{\tr})
  \widehat{\ve{P}}_t
  (\ve{I}-\ve{\kappa}_t\ve{k}_t^{\tr})^{\tr}
  +
  r_t\ve{\kappa}_t\ve{\kappa}_t^{\tr}.
  \label{eq:kam-candidate-cov}
\end{equation}
The Kalman write direction is the minimizer of total posterior uncertainty,
\begin{equation}
  \ve{\kappa}_t^\star
  =
  \arg\min_{\ve{\kappa}_t}
  \Trace\!\left(\ve{P}_t(\ve{\kappa}_t)\right).
  \label{eq:kam-trace-objective}
\end{equation}
Expanding the trace gives
\begin{equation}
  \Trace\!\left(\ve{P}_t(\ve{\kappa}_t)\right)
  =
  \Trace(\widehat{\ve{P}}_t)
  -2\ve{\kappa}_t^{\tr}\widehat{\ve{P}}_t\ve{k}_t
  +
  \big(r_t+\ve{k}_t^{\tr}\widehat{\ve{P}}_t\ve{k}_t\big)
  \ve{\kappa}_t^{\tr}\ve{\kappa}_t .
  \label{eq:kam-trace-expanded}
\end{equation}
Taking the gradient with respect to $\ve{\kappa}_t$ and setting it to zero yields
\begin{equation}
  \boxed{
  \ve{\kappa}_t
  =
  \frac{\widehat{\ve{P}}_t\ve{k}_t}
  {r_t+\ve{k}_t^{\tr}\widehat{\ve{P}}_t\ve{k}_t}
  }.
  \label{eq:kam-optimal-gain}
\end{equation}
Therefore the exact Kalman Associative Memory update is
\begin{equation}
\begin{aligned}
  \widehat{\ve{S}}_t
  &=\ve{D}_t\ve{S}_{t-1},
  &
  \widehat{\ve{P}}_t
  &=\ve{D}_t\ve{P}_{t-1}\ve{D}_t^{\tr}+\ve{\Omega}_t,
  \\
  \ve{\kappa}_t
  &=
  \frac{\widehat{\ve{P}}_t\ve{k}_t}
  {r_t+\ve{k}_t^{\tr}\widehat{\ve{P}}_t\ve{k}_t},
  &
  \ve{S}_t
  &=
  \widehat{\ve{S}}_t
  +
  \ve{\kappa}_t
  \big(\ve{v}_t-\widehat{\ve{S}}_t^{\tr}\ve{k}_t\big)^{\tr},
  \\
  \ve{P}_t
  &=
  (\ve{I}-\ve{\kappa}_t\ve{k}_t^{\tr})\widehat{\ve{P}}_t.
\end{aligned}
\label{eq:kam-optimal-update}
\end{equation}
This is the reference update. The remaining question is which part of it can be made
tractable and scan-computable without losing the link between uncertainty, transition, and
write strength.

\section{Dense covariance and dense information before approximation}

Before introducing any diagonal restriction, keep the covariance as a full dense matrix
$\ve{P}_t\in\R^{d_k\times d_k}$. The prediction step is exactly
\begin{equation}
  \widehat{\ve{P}}_t
  =
  \ve{D}_t\ve{P}_{t-1}\ve{D}_t^{\tr}+\ve{\Omega}_t.
  \label{eq:dense-predict}
\end{equation}
After the measurement at key $\ve{k}_t$, the Kalman gain and covariance update are
\begin{equation}
  \ve{\kappa}_t
  =
  \frac{\widehat{\ve{P}}_t\ve{k}_t}
  {r_t+\ve{k}_t^{\tr}\widehat{\ve{P}}_t\ve{k}_t},
  \qquad
  \ve{P}_t
  =
  (\ve{I}-\ve{\kappa}_t\ve{k}_t^{\tr})\widehat{\ve{P}}_t.
  \label{eq:dense-cov-update}
\end{equation}
Substituting the gain into the covariance update gives the dense posterior covariance
explicitly:
\begin{equation}
\begin{aligned}
  \ve{P}_t
  &=
  \widehat{\ve{P}}_t
  -
  \frac{\widehat{\ve{P}}_t\ve{k}_t}
       {r_t+\ve{k}_t^{\tr}\widehat{\ve{P}}_t\ve{k}_t}
  \ve{k}_t^{\tr}\widehat{\ve{P}}_t \\
  &=
  \widehat{\ve{P}}_t
  -
  \frac{
    \widehat{\ve{P}}_t\ve{k}_t\ve{k}_t^{\tr}\widehat{\ve{P}}_t
  }
  {r_t+\ve{k}_t^{\tr}\widehat{\ve{P}}_t\ve{k}_t}.
\end{aligned}
  \label{eq:dense-cov-expanded}
\end{equation}
This is still exact and dense.

Now define dense information matrices
\[
  \ve{C}_t=\ve{P}_t^{-1},
  \qquad
  \widehat{\ve{C}}_t=\widehat{\ve{P}}_t^{-1}.
\]
The rank-one Woodbury identity gives
\begin{equation}
  \left(\widehat{\ve{C}}_t+\frac{1}{r_t}\ve{k}_t\ve{k}_t^{\tr}\right)^{-1}
  =
  \widehat{\ve{P}}_t
  -
  \frac{
    \widehat{\ve{P}}_t\ve{k}_t\ve{k}_t^{\tr}\widehat{\ve{P}}_t
  }
  {r_t+\ve{k}_t^{\tr}\widehat{\ve{P}}_t\ve{k}_t}.
  \label{eq:dense-woodbury}
\end{equation}
The right-hand side is exactly \eqref{eq:dense-cov-expanded}. Therefore the same dense
measurement update in information form is
\begin{equation}
  \boxed{
  \ve{C}_t
  =
  \widehat{\ve{C}}_t+\frac{1}{r_t}\ve{k}_t\ve{k}_t^{\tr}.
  }
  \label{eq:exact-info-update}
\end{equation}
So if we allow dense $\ve{P}_t$ or dense $\ve{C}_t$, there is no approximation here. The
measurement adds a rank-one dense information matrix.

\paragraph{What becomes difficult.}
The dense route is mathematically clean but not the desired linear-attention primitive.
First, the state itself costs $O(d_k^2)$ per head instead of $O(d_k)$. Computing the gain
requires the dense matrix-vector product $\widehat{\ve{P}}_t\ve{k}_t$ and the quadratic form
$\ve{k}_t^{\tr}\widehat{\ve{P}}_t\ve{k}_t$, which are $O(d_k^2)$ per token. Second, the
covariance recursion is sequential: the gain at time $t$ depends on the dense posterior
uncertainty produced by all previous keys. Third, dense information form only makes the
measurement update simple. The prediction step becomes
\begin{equation}
  \widehat{\ve{C}}_t
  =
  \left(\ve{D}_t\ve{C}_{t-1}^{-1}\ve{D}_t^{\tr}+\ve{\Omega}_t\right)^{-1},
  \label{eq:dense-info-predict-hard}
\end{equation}
which requires an inverse when the process noise $\ve{\Omega}_t$ is additive. Thus dense
$\ve{C}_t$ is easy to update after a measurement but hard to predict through process noise,
while dense $\ve{P}_t$ is easy to predict but expensive to update and use for gains.

\paragraph{What is $\widehat{\ve{C}}_t$ used for?}
The Kalman gain itself is naturally written in terms of the predicted covariance:
\begin{equation}
  \ve{\kappa}_t
  =
  \frac{\widehat{\ve{P}}_t\ve{k}_t}
  {r_t+\ve{k}_t^{\tr}\widehat{\ve{P}}_t\ve{k}_t}.
  \label{eq:gain-needs-phat}
\end{equation}
Thus the gain does not need $\widehat{\ve{C}}_t$ as a matrix; it needs the vector
\begin{equation}
  \ve{x}_t=\widehat{\ve{P}}_t\ve{k}_t
  =
  \widehat{\ve{C}}_t^{-1}\ve{k}_t,
  \qquad
  \ve{k}_t^{\tr}\widehat{\ve{P}}_t\ve{k}_t=\ve{k}_t^{\tr}\ve{x}_t.
  \label{eq:gain-from-chat-solve}
\end{equation}
If $\widehat{\ve{C}}_t$ is dense, computing $\ve{x}_t$ means solving the linear system
$\widehat{\ve{C}}_t\ve{x}_t=\ve{k}_t$, which is not cheaper than dense covariance algebra in
general.

One can compute the gain ingredients directly from $\ve{C}_{t-1}$ without explicitly forming
$\widehat{\ve{C}}_t$:
\begin{equation}
\begin{aligned}
  \widehat{\ve{P}}_t\ve{k}_t
  &=
  \big(\ve{D}_t\ve{C}_{t-1}^{-1}\ve{D}_t^{\tr}+\ve{\Omega}_t\big)\ve{k}_t \\
  &=
  \ve{D}_t\ve{C}_{t-1}^{-1}(\ve{D}_t^{\tr}\ve{k}_t)
  +\ve{\Omega}_t\ve{k}_t.
\end{aligned}
  \label{eq:gain-direct-from-ctminus}
\end{equation}
This is exact if the solve with $\ve{C}_{t-1}$ is exact. It is useful conceptually because it
shows that $\widehat{\ve{C}}_t$ is mainly a predicted information state for the measurement
update and scan recursion, not a required intermediate for the gain. In the diagonal case,
\eqref{eq:gain-direct-from-ctminus} becomes coordinatewise:
\begin{equation}
  \widehat{p}_{t,i}k_{t,i}
  =
  \left(\frac{\alpha_{t,i}^2}{c_{t-1,i}}+\omega_{t,i}\right)k_{t,i},
  \qquad
  \ve{k}_t^{\tr}\widehat{\ve{P}}_t\ve{k}_t
  =
  \sum_i
  \left(\frac{\alpha_{t,i}^2}{c_{t-1,i}}+\omega_{t,i}\right)k_{t,i}^2.
  \label{eq:diag-gain-direct-from-cprev}
\end{equation}
Equivalently, $\widehat{c}_{t,i}=1/\widehat{p}_{t,i}$ is used when updating the information
state after the measurement.

\paragraph{Can we parameterize $\ve{\Omega}_t\ve{k}_t$ directly?}
The term $\ve{\Omega}_t\ve{k}_t$ is the process-noise contribution to the gain numerator.
Define
\[
  \ve{\xi}_t\triangleq \ve{\Omega}_t\ve{k}_t .
\]
Then the gain ingredients can be written as
\begin{equation}
\begin{aligned}
  \widehat{\ve{P}}_t\ve{k}_t
  &=
  \ve{D}_t\ve{C}_{t-1}^{-1}(\ve{D}_t^{\tr}\ve{k}_t)
  +\ve{\xi}_t,\\
  \ve{k}_t^{\tr}\widehat{\ve{P}}_t\ve{k}_t
  &=
  (\ve{D}_t^{\tr}\ve{k}_t)^{\tr}\ve{C}_{t-1}^{-1}
  (\ve{D}_t^{\tr}\ve{k}_t)
  +\ve{k}_t^{\tr}\ve{\xi}_t .
\end{aligned}
  \label{eq:gain-direct-xi}
\end{equation}
So for the gain alone, it is enough to know the vector action $\ve{\xi}_t$ and the scalar
$\ve{k}_t^{\tr}\ve{\xi}_t$; one does not need the full matrix $\ve{\Omega}_t$.

However, $\ve{\xi}_t$ is not automatically a valid process covariance. To remain
Kalman-consistent, there must exist a symmetric positive semidefinite $\ve{\Omega}_t$ such that
$\ve{\xi}_t=\ve{\Omega}_t\ve{k}_t$ and
$\ve{k}_t^{\tr}\ve{\xi}_t=\ve{k}_t^{\tr}\ve{\Omega}_t\ve{k}_t\ge 0$. In the diagonal case this
means
\[
  \ve{\xi}_t=\ve{\omega}_t\odot\ve{k}_t,
  \qquad
  \ve{\omega}_t\ge 0,
  \qquad
  \ve{k}_t^{\tr}\ve{\xi}_t=\sum_i\omega_{t,i}k_{t,i}^2.
\]
This is equivalent to parameterizing the diagonal process noise $\ve{\omega}_t$; it does not
change the uncertainty scan. If instead we parameterize an arbitrary vector $\ve{\xi}_t$ and a
positive scalar denominator contribution directly, the gain can be made easier to compute and
remains input-only once the prefix information state is known. But then the update is no
longer the exact Kalman covariance recursion unless we also define a compatible
$\ve{\Omega}_t$ and update the uncertainty state with it. In that case it should be viewed as a
learned gain approximation rather than a full process-noise model.

\paragraph{Options from the dense information prediction.}
Strictly, \eqref{eq:dense-info-predict-hard} defines the predicted information
$\widehat{\ve{C}}_t$, not the posterior information $\ve{C}_t$. From this equation, the
available paths are:
\begin{enumerate}[leftmargin=1.6em,itemsep=3pt]
  \item \textbf{Stay exact and dense.} Compute
  $\widehat{\ve{C}}_t=(\ve{D}_t\ve{C}_{t-1}^{-1}\ve{D}_t^{\tr}+\ve{\Omega}_t)^{-1}$ directly.
  This keeps the Kalman filter exact, but requires dense inverses or dense linear solves.

  \item \textbf{Use covariance for prediction.} Keep the uncertainty state in covariance form
  rather than information form. Then the process step is immediate:
  \[
    \widehat{\ve{P}}_t
    =
    \ve{D}_t\ve{P}_{t-1}\ve{D}_t^{\tr}+\ve{\Omega}_t .
  \]
  This avoids the difficult inverse in \eqref{eq:dense-info-predict-hard}. The gain is then
  computed from this predicted covariance:
  \[
    \ve{\kappa}_t
    =
    \frac{\widehat{\ve{P}}_t\ve{k}_t}
    {r_t+\ve{k}_t^{\tr}\widehat{\ve{P}}_t\ve{k}_t}.
  \]
  After the measurement, the covariance state for the next token is
  \[
    \ve{P}_t
    =
    (\ve{I}-\ve{\kappa}_t\ve{k}_t^{\tr})\widehat{\ve{P}}_t.
  \]
  Therefore the recurrence is
  \[
    \ve{P}_{t-1}
    \longrightarrow
    \widehat{\ve{P}}_t
    \longrightarrow
    \ve{\kappa}_t
    \longrightarrow
    \ve{P}_t .
  \]
  If implemented literally, this recurrence is sequential. If we ignore storage and compute
  cost, however, the dense covariance recursion is parallelizable in principle. Each token
  defines a Riccati map $\ve{P}_t=F_t(\ve{P}_{t-1})$, and prefix covariances are
  \[
    \ve{P}_t=(F_t\circ F_{t-1}\circ\cdots\circ F_1)(\ve{P}_0).
  \]
  Function composition is associative. More concretely, when $\ve{D}_t$ is invertible, the
  prediction and measurement updates are matrix fractional maps. The prediction can be
  written as
  \[
    \widehat{\ve{P}}_t
    =
    (\ve{D}_t\ve{P}_{t-1}+\ve{\Omega}_t\ve{D}_t^{-\tr})
    (\ve{D}_t^{-\tr})^{-1},
  \]
  and the measurement update can be written as
  \[
    \ve{P}_t
    =
    \widehat{\ve{P}}_t
    \left(\ve{I}+\frac{1}{r_t}\ve{k}_t\ve{k}_t^{\tr}\widehat{\ve{P}}_t\right)^{-1}.
  \]
  These fractional maps compose by multiplying their block-matrix representations, so the
  dense uncertainty prefix can be computed by a scan in principle. The issue is therefore not
  mathematical associativity. The issue is that this is a dense Riccati scan, not the cheap
  affine memory scan used by linear attention; after the scan one still has to recover
  $\widehat{\ve{P}}_t$ and form $\widehat{\ve{P}}_t\ve{k}_t$ for every token to get the gains.
  Converting to information would mean forming
  $\widehat{\ve{C}}_t=\widehat{\ve{P}}_t^{-1}$, which is another dense inverse, so this path
  usually stays in covariance form.

  The memory cost is the main practical obstacle. At autoregressive inference time, a dense
  covariance filter stores one $\ve{P}_t\in\R^{d_k\times d_k}$ per head, i.e.\ $O(d_k^2)$
  uncertainty state instead of the $O(d_k)$ state of a diagonal tracker. In parallel training,
  a dense Riccati scan must materialize either dense prefix covariances or dense fractional-map
  scan elements for many tokens. Each token's fractional map is represented by a
  $2d_k\times 2d_k$ block matrix, equivalently four dense $d_k\times d_k$ blocks, so the
  scan state is still $O(d_k^2)$ per token. Across a length-$T$ sequence this is
  $O(Td_k^2)$ per head just for the uncertainty scan, before storing the usual key, value,
  gain, and memory-scan quantities. The diagonal approximation reduces this uncertainty-scan
  storage to $O(Td_k)$.

  \item \textbf{Remove additive process noise.} If $\ve{\Omega}_t=0$ and $\ve{D}_t$ is
  invertible, then the information prediction is exact and simple:
  \[
    \widehat{\ve{C}}_t=\ve{D}_t^{-\tr}\ve{C}_{t-1}\ve{D}_t^{-1}.
  \]
  This is tractable algebraically, but it removes the additive uncertainty growth that
  represents drift.

  \item \textbf{Restrict to a diagonal family.} If $\ve{D}_t$, $\ve{\Omega}_t$, and
  $\ve{C}_{t-1}$ are diagonal, the inverse is coordinatewise:
  \[
    \widehat{c}_{t,i}
    =
    \frac{c_{t-1,i}}{\alpha_{t,i}^2+\omega_{t,i}c_{t-1,i}}.
  \]
  This is the tractable path used below, but the following measurement update must be
  projected back to diagonal.

  \item \textbf{Use structured covariance.} Keep blocks, diagonal-plus-low-rank factors, or
  chunkwise dense states. These approximate the dense equation more closely than a diagonal
  vector, but their scan and gain computations are more expensive.
\end{enumerate}

This is the point where tractability enters. The diagonal state below is not needed for the
Kalman derivation; it is a computational restriction that keeps the uncertainty state
linear in $d_k$ and makes a scan-compatible approximation possible.

\section{Diagonal information as the tractable state}

Because the dense recursions above are too expensive for a linear-attention layer, we
restrict the uncertainty state to the space of diagonal covariance matrices, the simplest
family that is still more expressive than a scalar gain:
\begin{equation}
  \ve{P}_t\in\{\diag(\ve{p}):\ve{p}\in\R_+^{d_k}\},
  \qquad
  \ve{P}_t=\diag(\ve{p}_t),
  \qquad
  \ve{C}_t=\ve{P}_t^{-1}=\diag(\ve{c}_t).
  \label{eq:diag-cov-family}
\end{equation}
Here $\ve{p}_t$ is per-channel uncertainty and $\ve{c}_t$ is per-channel information.
The information parameterization $\ve{C}_t$ is useful, but it is only one possible way to
move forward; the exact object under the diagonal restriction is still the covariance
prediction below.

\subsection{Transition-aware diagonal prediction}

With a diagonal process model $\ve{D}_t=\diag(\ve{\alpha}_t)$ and diagonal process noise
$\ve{\Omega}_t=\diag(\ve{\omega}_t)$, the Kalman covariance prediction remains pointwise:
\begin{equation}
  \widehat{\ve{p}}_t
  =
  \ve{\alpha}_t^2\odot\ve{p}_{t-1}+\ve{\omega}_t .
  \tag{21}
  \label{eq:diag-predict-vector}
\end{equation}
The gain is then the diagonal approximation of the Kalman gain,
\begin{equation}
  \ve{\kappa}_t
  =
  \frac{\widehat{\ve{p}}_t\odot\ve{k}_t}
  {r_t+\sum_i\widehat{p}_{t,i}k_{t,i}^2}.
  \tag{22}
  \label{eq:diag-gain-vector}
\end{equation}
Unlike the fixed-gain delta rule, in the exact prediction \eqref{eq:diag-predict-vector} the
gain $\ve{\kappa}_t$ depends on the same process model that predicts the memory: if a channel
is forgotten or assigned high process noise, its uncertainty changes before the write is
formed. How much of this coupling survives once the recursion is made scan-computable
depends on the parameterization.

Per channel, the covariance predict is
\begin{equation}
  \boxed{\widehat{p}_t=\alpha_t^2p_{t-1}+\omega_t.}
  \tag{67}
  \label{eq:diag-predict-scalar}
\end{equation}
This is the diagonal form of \eqref{eq:diag-predict-vector}. It is the only exact predicted
variance; every expression below is \eqref{eq:diag-predict-scalar} written in different
variables or followed by an approximation.

\subsection*{Options after the exact predicted variance}

At this point the derivation is still exact under the diagonal covariance restriction. There
are several ways to go forward:
\begin{enumerate}[leftmargin=1.6em,itemsep=3pt]
  \item \textbf{Keep covariance variables.} Track $\ve{p}_t$ directly and compute
  $\widehat{\ve{p}}_t$ by \eqref{eq:diag-predict-vector}. This is closest to the Kalman gain,
  but the posterior covariance update after measurement is not diagonal unless we project.

  \item \textbf{Use information variables.} Track $\ve{c}_t=\ve{p}_t^{-1}$. This makes the
  Gaussian measurement update additive in information form and can lead to a scan-computable
  recursion. This is one option, not the only one.

  \item \textbf{Project to an isotropic scalar.} Replace $\ve{p}_t$ by $b_t\ve{1}$ or
  $\ve{P}_t$ by $b_t\ve{I}$. This preserves the delta-rule write direction
  $\ve{\kappa}_t\propto\ve{k}_t$, but loses per-channel uncertainty.

  \item \textbf{Use a richer structured family.} Track block-diagonal, diagonal-plus-low-rank,
  or chunk-boundary covariance states. These keep more correlations than a diagonal vector,
  but they are more expensive and require a different scan design.

  \item \textbf{Freeze or approximate the gain.} Compute gains from input-only or chunkwise
  prefix statistics and then use the standard affine memory scan. This is an implementation
  choice, not an exact Kalman covariance recursion.
\end{enumerate}

\section{Measurement boundary after choosing a diagonal state}

From the dense derivation above, the measurement update in information variables is
\eqref{eq:exact-info-update}. Now impose the diagonal state restriction. If
$\widehat{\ve{C}}_t=\diag(\widehat{\ve{c}}_t)$, then the measurement still gives
\begin{equation}
  \ve{C}_t
  =
  \diag(\widehat{\ve{c}}_t)+\frac{1}{r_t}\ve{k}_t\ve{k}_t^{\tr}.
  \label{eq:diag-prior-dense-posterior}
\end{equation}
The rank-one term is dense for a general key $\ve{k}_t$, so the posterior information matrix
is generally not diagonal:
\[
  \ve{C}_t\notin
  \{\diag(\ve{c}):\ve{c}\in\R_+^{d_k}\}.
\]
A vector state $\ve{c}_t$ cannot represent this posterior exactly. This is where exactness
stops for the diagonal family: any diagonal update after \eqref{eq:diag-prior-dense-posterior}
is a projection or approximation.

\subsection*{Options at the measurement boundary}

After reaching the dense posterior in \eqref{eq:diag-prior-dense-posterior}, the available
choices are:
\begin{enumerate}[leftmargin=1.6em,itemsep=3pt]
  \item \textbf{Keep the exact dense posterior.} This remains Kalman-exact, but it loses the
  diagonal state and is not the intended tractable linear-attention update.

  \item \textbf{Literal diagonal projection.} Keep only the diagonal of
  $\ve{k}_t\ve{k}_t^{\tr}/r_t$, giving
  $\ve{c}_t=\widehat{\ve{c}}_t+(\ve{k}_t\odot\ve{k}_t)/r_t$. This is the literal diagonal of the
  exact Fisher term, but it is poorly scaled for normalized high-dimensional keys.

  \item \textbf{Calibrated diagonal projection.} Standardize each key coordinate before taking
  the diagonal Fisher increment, giving the dimension-stable update derived below.

  \item \textbf{Isotropic projection.} Average the calibrated diagonal information into one
  scalar. This keeps a scalar uncertainty and recovers a gain direction proportional to
  $\ve{k}_t$.

  \item \textbf{Structured projection.} Keep block-diagonal or low-rank pieces of the
  rank-one Fisher term. This is closer to the exact posterior but more expensive than a
  diagonal vector.
\end{enumerate}

\section{Interlude: why softmax attention needs the $1/\sqrt{d_k}$ scale}
\label{sec:attention-sqrt-dk}

The calibration in the next section standardizes a key coordinate before using it. That is the
same operation that produces the familiar $1/\sqrt{d_k}$ factor in softmax attention, so it is
worth deriving there first, where the failure mode is easy to quantify. The conclusion of this
section is that $\sqrt{d_k}$ is not a tuning constant: it is the unique power of $d_k$ that makes
the attention logit an \emph{intensive} quantity, one whose distribution does not change as the
head dimension grows.

\subsection*{The logit is extensive in $d_k$}

Softmax attention scores a query $\ve{q}\in\R^{d_k}$ against keys
$\ve{k}_1,\ldots,\ve{k}_T\in\R^{d_k}$ through the inner product
\[
  s_j=\ve{q}^{\tr}\ve{k}_j=\sum_{i=1}^{d_k}q_ik_{j,i},
\]
and normalizes $\ve{s}$ with a softmax. At initialization the coordinates of $\ve{q}$ and
$\ve{k}_j$ are well modeled as independent, zero-mean, unit-variance: this is what a normalization
layer followed by a standard random projection produces. Then $\mathbb{E}[q_ik_{j,i}]=0$ and
$\operatorname{Var}(q_ik_{j,i})=\mathbb{E}[q_i^2]\mathbb{E}[k_{j,i}^2]=1$, so summing $d_k$
independent terms gives
\begin{equation}
  \mathbb{E}[s_j]=0,
  \qquad
  \operatorname{Var}(s_j)=d_k,
  \qquad
  s_j=\Theta(\sqrt{d_k}).
  \label{eq:logit-variance}
\end{equation}
The logit is therefore \emph{extensive}: it grows with the representation dimension. Nothing about
the modeling problem changed when $d_k$ grew from $64$ to $128$, but the scale of the quantity fed
to the softmax did.

\subsection*{What goes wrong: saturation and a vanishing softmax Jacobian}

The softmax is not scale-free. Writing $\ve{p}=\mathrm{softmax}(\ve{s})$, its Jacobian is
\begin{equation}
  \frac{\partial p_a}{\partial s_b}=p_a(\delta_{ab}-p_b),
  \label{eq:softmax-jacobian}
\end{equation}
which vanishes in both saturated limits: as $\ve{p}$ approaches a one-hot vector, every entry of
\eqref{eq:softmax-jacobian} tends to zero. Large logit spread therefore does not merely make the
attention pattern peaked at initialization, it removes the gradient that would let the pattern
change.

The two-key case makes the rate explicit. For distinct keys the gap $\Delta=s_1-s_2$ has
conditional variance $2\|\ve{q}\|^2\approx 2d_k$, so $|\Delta|=\Theta(\sqrt{2d_k})$, and the
gradient scale through the softmax is
\begin{equation}
  p(1-p)=\sigma(\Delta)\sigma(-\Delta)\le e^{-|\Delta|},
  \qquad
  \sigma(x)=\frac{1}{1+e^{-x}}.
  \label{eq:two-key-gradient}
\end{equation}
The suppression is exponential in $\sqrt{d_k}$:

\begin{center}
\begin{tabular}{rccc}
\hline
 & \multicolumn{2}{c}{unscaled $\ve{q}^{\tr}\ve{k}$} & scaled $\ve{q}^{\tr}\ve{k}/\sqrt{d_k}$ \\
$d_k$ & typical $|\Delta|$ & $p(1-p)$ & typical $|\Delta|$, $p(1-p)$ \\
\hline
$16$  & $5.7$  & $3.5\times10^{-3}$ & $1.41$, $\;0.158$ \\
$64$  & $11.3$ & $1.2\times10^{-5}$ & $1.41$, $\;0.158$ \\
$128$ & $16.0$ & $1.1\times10^{-7}$ & $1.41$, $\;0.158$ \\
\hline
\end{tabular}
\end{center}

\noindent
The scaled column is constant in $d_k$; that invariance is the entire point. With $T$ keys the same
effect appears as attention-entropy collapse: for logits of standard deviation $\sigma$ the largest
gap grows like $\sigma\sqrt{2\log T}$, so the attention distribution concentrates on one key as soon
as $\sigma\gg 1/\sqrt{2\log T}$.

\subsection*{Why the exponent is exactly $1/2$}

Consider the family $s^{(\gamma)}=d_k^{-\gamma}\,\ve{q}^{\tr}\ve{k}$. From
\eqref{eq:logit-variance},
\[
  \operatorname{Var}\!\left(s^{(\gamma)}\right)=d_k^{1-2\gamma}.
\]
Three regimes follow, and only one is usable:
\begin{enumerate}[leftmargin=1.6em,itemsep=3pt]
  \item $\gamma<\tfrac12$: the variance diverges with $d_k$. The softmax saturates, the Jacobian
  \eqref{eq:softmax-jacobian} collapses, and wider heads train worse than narrow ones.

  \item $\gamma>\tfrac12$: the variance vanishes with $d_k$. The softmax tends to the uniform
  distribution, attention degenerates to unweighted mean pooling, and the gradient reaching
  $\ve{q}$ and $\ve{k}$ is suppressed by the same factor.

  \item $\gamma=\tfrac12$: the variance is $1$ for every $d_k$. The logit distribution, the
  effective softmax temperature, and hence the usable learning rate are all independent of head
  dimension.
\end{enumerate}
So $1/\sqrt{d_k}$ is the unique scale under which one hyperparameter setting transfers across head
widths. Dividing by $d_k$ is not a safe ``more conservative'' choice; it fails in the opposite
direction.

\subsection*{The same statement as coordinate standardization}

Dividing the score by $\sqrt{d_k}$ can be moved onto the key. Under $\ell_2$ normalization of
$\ve{q}$ and $\ve{k}$ --- the convention used by the mixers in this work --- the raw score is
$\cos\theta$, whose variance for independent directions is $1/d_k$ rather than $d_k$. The
standardizing correction is then a \emph{multiplication}:
\begin{equation}
  \sqrt{d_k}\,\ve{q}^{\tr}\ve{k}
  =\ve{q}^{\tr}\widetilde{\ve{k}},
  \qquad
  \widetilde{\ve{k}}=\sqrt{d_k}\,\ve{k},
  \qquad
  \mathbb{E}\big[\widetilde k_i^{\,2}\big]\approx 1.
  \label{eq:standardized-key}
\end{equation}
Both conventions express one requirement: standardize the per-coordinate contribution so that the
sum over $d_k$ coordinates is $O(1)$. The unnormalized convention divides the sum; the normalized
convention rescales the coordinate. The object $\widetilde{\ve{k}}=\sqrt{d_k}\,\ve{k}$ in
\eqref{eq:standardized-key} is exactly the standardized coordinate used in the next section.

This also explains why the calibrated filter carries $d_k$ where attention carries $\sqrt{d_k}$.
The attention score is \emph{linear} in the key, so standardization contributes one factor of
$\sqrt{d_k}$. The Fisher information $\widetilde{\ve{k}}\widetilde{\ve{k}}^{\tr}/r_t$ is
\emph{quadratic} in the key, so the same standardization contributes two, giving the per-channel
increment $d_kk_{t,i}^2/r_t$ of \eqref{eq:calibrated-diag-update}. The two constants are the same
argument applied to a linear and a quadratic functional of the same standardized key.

\subsection*{Scope of the argument}

Equation~\eqref{eq:logit-variance} assumes $\ve{q}$ and $\ve{k}$ are independent, which holds at
initialization but not after training: both are learned functions of the same input, and an
attention head that has learned to retrieve produces aligned pairs. For aligned $\ve{q},\ve{k}$ the
inner product is $\Theta(\|\ve{q}\|\|\ve{k}\|)=\Theta(d_k)$, and the $1/\sqrt{d_k}$ scale still
leaves a residual $\Theta(\sqrt{d_k})$ growth. This is the known origin of logit growth and
attention-entropy collapse late in training, and the reason two later corrections exist: the
maximal-update parameterization ($\mu$P) scales attention logits by $1/d_k$ rather than
$1/\sqrt{d_k}$ precisely because it assumes the correlated regime, and QK-normalization projects
$\ve{q},\ve{k}$ to the sphere and reintroduces an explicit learned temperature. The
$1/\sqrt{d_k}$ factor should therefore be read as the correct \emph{initialization-time}
standardization, not as a scale invariant that training preserves.

\section{Calibrated diagonal information update}

If we choose a diagonal vector state, we need to replace the dense Fisher matrix
$\ve{k}_t\ve{k}_t^{\tr}/r_t$ by a diagonal increment. The literal diagonal increment is
\[
  \frac{1}{r_t}\ve{k}_t\odot\ve{k}_t.
\]
For $\ell_2$-normalized keys with spread-out energy, a typical coordinate satisfies
$\mathbb{E}[k_{t,i}^2]\approx 1/d_k$. Therefore the literal per-channel increment has size
\[
  \mathbb{E}\!\left[\frac{k_{t,i}^2}{r_t}\right]
  \approx
  \frac{1}{d_k r_t},
\]
which vanishes as the key dimension grows. That scaling makes the diagonal tracker depend on
the arbitrary representation dimension rather than on the observation reliability $r_t$.

To get a dimension-stable per-channel Fisher increment, use the standardized coordinate
\[
  \widetilde{k}_{t,i}=\sqrt{d_k}\,k_{t,i}.
\]
For normalized spread-out keys,
\[
  \mathbb{E}[\widetilde{k}_{t,i}^{\,2}]
  =
  d_k\mathbb{E}[k_{t,i}^2]
  \approx 1.
\]
The scalar Fisher information contributed by coordinate $i$ is therefore
\[
  \frac{\widetilde{k}_{t,i}^{\,2}}{r_t}
  =
  \frac{d_k}{r_t}k_{t,i}^2.
\]
This gives the calibrated diagonal update
\begin{equation}
  \boxed{
  \ve{c}_t
  =
  \widehat{\ve{c}}_t
  +
  \frac{d_k}{r_t}\,\ve{k}_t\odot\ve{k}_t.
  }
  \label{eq:calibrated-diag-update}
\end{equation}
Equivalently, per channel,
\begin{equation}
  c_{t,i}
  =
  \widehat{c}_{t,i}
  +
  \frac{d_k}{r_t}k_{t,i}^2.
  \label{eq:calibrated-diag-update-coordinate}
\end{equation}

\subsection*{Where approximation enters}

The prediction
\[
  \widehat{p}_{t,i}=\alpha_{t,i}^2p_{t-1,i}+\omega_{t,i}
\]
is exact under the diagonal covariance restriction. The dense information update
\[
  \ve{C}_t
  =
  \widehat{\ve{C}}_t+\ve{k}_t\ve{k}_t^{\tr}/r_t
\]
is exact for the linear--Gaussian measurement. The calibrated diagonal update
\eqref{eq:calibrated-diag-update} is not the exact posterior; it is a dimension-stable
diagonal projection of the measurement information. From this point onward, any
scan-computable diagonal recursion, including an information-form or M\"obius recurrence,
is an approximation to the exact dense Kalman posterior.

\subsection*{Options after the calibrated projection}

Once the calibrated diagonal update is chosen, the remaining design options are:
\begin{enumerate}[leftmargin=1.6em,itemsep=3pt]
  \item \textbf{Free additive process noise.} Keep $\omega_{t,i}$ as an input-only additive
  variance. In information variables this produces a per-channel linear-fractional
  recurrence that can be scanned with $2\times2$ matrices.

  \item \textbf{State-proportional process noise.} Choose $\omega_t\propto p_{t-1}$ so the
  information recursion becomes affine. This is cheaper, but it restricts the process-noise
  model.

  \item \textbf{Isotropic information.} Average the diagonal information into one scalar if the
  goal is to preserve the delta-rule write direction exactly.

  \item \textbf{Richer covariance families.} Move beyond the diagonal projection with blocks,
  low-rank terms, or chunkwise dense updates if the diagonal approximation is too weak.
\end{enumerate}

\section*{Note: language-modeling objective}

The trace objective in \eqref{eq:kam-trace-objective} is a local Bayesian criterion for the
associative-memory update. Under the linear--Gaussian model, it chooses the Kalman/MMSE write
direction by minimizing posterior covariance. This is a useful adaptive-gain prior, but it is
not the objective optimized by an autoregressive language model. A language model is trained
to minimize future next-token negative log likelihood under a finite memory budget.

For the language-modeling use case, the more aligned local question is therefore:
\begin{quote}
  \emph{Choose the write gain that most improves future next-token prediction under limited
  memory capacity.}
\end{quote}
Writing
\[
  \ve{S}_t(\ve{\kappa})
  =
  \widehat{\ve{S}}_t
  +
  \ve{\kappa}
  \big(\ve{v}_t-\widehat{\ve{S}}_t^{\tr}\ve{k}_t\big)^{\tr},
\]
a direct idealization is
\begin{equation}
  \ve{\kappa}_t^\star
  \approx
  \arg\min_{\ve{\kappa}}
  \mathbb{E}_{1\le h\le H}
  \left[
    \ell_{\mathrm{LM}}\!\left(
      x_{t+h}\mid x_{<t+h},\ve{S}_t(\ve{\kappa})
    \right)
  \right]
  +
  \lambda_{\mathrm{int}}\,
  \mathcal{I}_t(\ve{\kappa}),
  \label{eq:lm-predictive-memory-objective}
\end{equation}
where $\mathcal{I}_t$ measures interference with useful contents already stored in the
memory. This objective is not meant to be solved literally inside the layer. It clarifies
what the Kalman update should approximate if it is to be useful for language modeling.

\paragraph{Short-horizon predictive loss.}
The most direct proxy is to prefer writes that reduce cross entropy over a short future
horizon:
\[
  \ve{\kappa}_t
  \quad\text{should reduce}\quad
  \ell_{\mathrm{LM}}(x_{t+1:t+H}\mid x_{\le t},\ve{S}_t).
\]
This says that a good write is one that helps future prediction, not merely one that reduces
geometric uncertainty.

\paragraph{Future-read reconstruction.}
The memory should store information that future queries will actually read. A corresponding
proxy is
\[
  \mathbb{E}_{\tau>t}
  \left[
    \left\|
      \ve{S}_t(\ve{\kappa})^{\tr}\ve{q}_{\tau}
      -
      \ve{y}_{\tau}
    \right\|_2^2
  \right],
\]
where $\ve{q}_{\tau}$ denotes a future read query and $\ve{y}_{\tau}$ denotes the target
content for that query. This differs from $\Trace(\ve{P}_t)$ because language modeling does
not value all key-space directions equally; it values directions that future tokens will
query.

\paragraph{Surprise-weighted writing.}
A language-model memory should write more when the current token contributes useful new
information. Thus the gain should depend not only on key-space uncertainty, but also on a
causal surprise signal:
\[
  \text{write strength}
  \propto
  \text{useful surprise}
  \times
  \text{retrievability}
  \times
  \text{non-redundancy}.
\]
The associative-memory innovation is one local, non-leaking proxy:
\[
  \ve{\delta}_t
  =
  \ve{v}_t-\widehat{\ve{S}}_t^{\tr}\ve{k}_t,
  \qquad
  s_t
  =
  \operatorname{sg}\!\left(
    \frac{1}{d_v}\|\ve{\delta}_t\|_2^2
  \right),
\]
where $\operatorname{sg}(\cdot)$ denotes stop-gradient. A conservative implementation can
let $s_t$ modulate the Kalman confidence parameters:
\[
  r_t
  =
  r_{\min}+\operatorname{softplus}(\rho_t-a\,\widetilde{s}_t),
  \qquad
  \omega_{t,i}
  =
  \omega_{\min}+\operatorname{softplus}(\eta_{t,i}+b_i\,\widetilde{s}_t).
\]
Large innovation can then decrease $r_t$ or increase $\omega_t$, allowing a larger update.
The modulation coefficients should be initialized near zero and used with floors or clamps,
so the model starts as ordinary KLA and learns whether surprise-dependent plasticity helps.

\paragraph{Capacity-aware utility.}
Because the recurrent state is small relative to the sequence, every write competes with
existing contents. A language-model-aligned update should therefore optimize a net utility:
\[
  \text{future CE improvement}
  -
  \text{interference cost}.
\]
This is the most relevant framing for KLA and DeltaNet-style memories. A write is useful
only if it improves future recall or prediction more than it harms information already stored
for future reads.

In this framing, Kalman uncertainty minimization provides a principled adaptive-gain prior.
For language models, the actual selection pressure is predictive memory utility: writes
should reduce future language-model loss under capacity and interference constraints.
Signals derived from top-level entropy or realized negative log likelihood must be used
causally, for example to modulate the next update, so that the memory update does not leak
the target token into the current hidden state.

\end{document}
